\documentclass[11pt]{scrreprt}
%%fakesection Load packages
% These are evan.sty
\usepackage{amsmath,amssymb,amsthm}
\usepackage{mathrsfs}
\usepackage[usenames,svgnames,dvipsnames]{xcolor}
\usepackage{hyperref}
\usepackage[nameinlink]{cleveref}
\usepackage{textcomp}
\usepackage{enumerate}
\usepackage[obeyFinal,textsize=scriptsize,shadow]{todonotes}
\usepackage{mathtools}
\usepackage{microtype}

\usepackage[normalem]{ulem}
\usepackage{stmaryrd}
\usepackage{wasysym}
\usepackage{multirow}
\usepackage{prerex}

%%fakesection evan.sty macros
%Small commands
\newcommand{\cbrt}[1]{\sqrt[3]{#1}}
\newcommand{\floor}[1]{\left\lfloor #1 \right\rfloor}
\newcommand{\ceiling}[1]{\left\lceil #1 \right\rceil}
\newcommand{\mailto}[1]{\href{mailto:#1}{\texttt{#1}}}
\newcommand{\ol}{\overline}
\newcommand{\ul}{\underline}
\newcommand{\wt}{\widetilde}
\newcommand{\wh}{\widehat}
\newcommand{\eps}{\varepsilon}
%\renewcommand{\iff}{\Leftrightarrow}
%\renewcommand{\implies}{\Rightarrow}
\newcommand{\vocab}[1]{\textbf{\color{blue} #1}}
\newcommand{\catname}{\mathsf}
\newcommand{\hrulebar}{
  \par\hspace{\fill}\rule{0.95\linewidth}{.7pt}\hspace{\fill}
  \par\nointerlineskip \vspace{\baselineskip}
}
\newcommand{\half}{\frac{1}{2}}

%For use in author command
\newcommand{\plusemail}[1]{\\ \normalfont \texttt{\mailto{#1}}}

%More commands and math operators
\DeclareMathOperator{\cis}{cis}
\DeclareMathOperator*{\lcm}{lcm}
\DeclareMathOperator*{\argmin}{arg min}
\DeclareMathOperator*{\argmax}{arg max}

%Convenient Environments
\newenvironment{soln}{\begin{proof}[Solution]}{\end{proof}}
\newenvironment{parlist}{\begin{inparaenum}[(i)]}{\end{inparaenum}}
\newenvironment{gobble}{\setbox\z@\vbox\bgroup}{\egroup}

%Inequalities
\newcommand{\cycsum}{\sum_{\mathrm{cyc}}}
\newcommand{\symsum}{\sum_{\mathrm{sym}}}
\newcommand{\cycprod}{\prod_{\mathrm{cyc}}}
\newcommand{\symprod}{\prod_{\mathrm{sym}}}

%From H113 "Introduction to Abstract Algebra" at UC Berkeley
\newcommand{\CC}{\mathbb C}
\newcommand{\FF}{\mathbb F}
\newcommand{\NN}{\mathbb N}
\newcommand{\QQ}{\mathbb Q}
\newcommand{\RR}{\mathbb R}
\newcommand{\ZZ}{\mathbb Z}
\newcommand{\charin}{\text{ char }}
\DeclareMathOperator{\sign}{sign}
\DeclareMathOperator{\Aut}{Aut}
\DeclareMathOperator{\Inn}{Inn}
\DeclareMathOperator{\Syl}{Syl}
\DeclareMathOperator{\Gal}{Gal}
\DeclareMathOperator{\GL}{GL} % General linear group
\DeclareMathOperator{\SL}{SL} % Special linear group

%From Kiran Kedlaya's "Geometry Unbound"
\newcommand{\abs}[1]{\left\lvert #1 \right\rvert}
\newcommand{\norm}[1]{\left\lVert #1 \right\rVert}
\newcommand{\dang}{\measuredangle} %% Directed angle
\newcommand{\ray}[1]{\overrightarrow{#1}} 
\newcommand{\seg}[1]{\overline{#1}}
\newcommand{\arc}[1]{\wideparen{#1}}

%From M275 "Topology" at SJSU
\newcommand{\id}{\mathrm{id}}
\newcommand{\taking}[1]{\xrightarrow{#1}}
\newcommand{\inv}{^{-1}}

%From M170 "Introduction to Graph Theory" at SJSU
\DeclareMathOperator{\diam}{diam}
\DeclareMathOperator{\ord}{ord}
\newcommand{\defeq}{\overset{\mathrm{def}}{=}}

%From the USAMO .tex filse
\newcommand{\st}{^{\text{st}}}
\newcommand{\nd}{^{\text{nd}}}
\renewcommand{\th}{^{\text{th}}}
\newcommand{\dg}{^\circ}
\newcommand{\ii}{\item}

% From Math 55 and Math 145 at Harvard
\newenvironment{subproof}[1][Proof]{%
\begin{proof}[#1] \renewcommand{\qedsymbol}{$\blacksquare$}}%
{\end{proof}}

\newcommand{\liff}{\leftrightarrow}
\newcommand{\lthen}{\rightarrow}
\newcommand{\opname}{\operatorname}
\newcommand{\surjto}{\twoheadrightarrow}
\newcommand{\injto}{\hookrightarrow}
\newcommand{\On}{\mathrm{On}} % ordinals
\DeclareMathOperator{\img}{im} % Image
\DeclareMathOperator{\Img}{Im} % Image
\DeclareMathOperator{\coker}{coker} % Cokernel
\DeclareMathOperator{\Coker}{Coker} % Cokernel
\DeclareMathOperator{\Ker}{Ker} % Kernel
\DeclareMathOperator{\rank}{rank}
\DeclareMathOperator{\Spec}{Spec} % spectrum
\DeclareMathOperator{\Tr}{Tr} % trace
\DeclareMathOperator{\pr}{pr} % projection
\DeclareMathOperator{\ext}{ext} % extension
\DeclareMathOperator{\pred}{pred} % predecessor
\DeclareMathOperator{\dom}{dom} % domain
\DeclareMathOperator{\ran}{ran} % range
\DeclareMathOperator{\Hom}{Hom} % homomorphism
\DeclareMathOperator{\End}{End} % endomorphism

% Things Lie
\newcommand{\kb}{\mathfrak b}
\newcommand{\kg}{\mathfrak g}
\newcommand{\kh}{\mathfrak h}
\newcommand{\kn}{\mathfrak n}
\newcommand{\ku}{\mathfrak u}
\newcommand{\kz}{\mathfrak z}
\DeclareMathOperator{\Ext}{Ext} % Ext functor
\DeclareMathOperator{\Tor}{Tor} % Tor functor
\newcommand{\gl}{\opname{\mathfrak{gl}}} % frak gl group
\renewcommand{\sl}{\opname{\mathfrak{sl}}} % frak gl group

% More script letters etc.
\newcommand{\SA}{\mathscr A}
\newcommand{\SB}{\mathscr B}
\newcommand{\SC}{\mathscr C}
\newcommand{\SD}{\mathscr D}
\newcommand{\SF}{\mathscr F}
\newcommand{\SG}{\mathscr G}
\newcommand{\SH}{\mathscr H}
\newcommand{\OO}{\mathcal O}

%% Napkin commands
\newcommand{\prototype}[1]{
	\emph{{\color{red} Prototypical example for this section:} #1} \par\medskip
}
\newenvironment{moral}{%
	\begin{mdframed}[linecolor=green!70!black]%
	\bfseries\color{green!70!black}}%
	{\end{mdframed}}

%%% END EXCERPT OF EVAN.STY
%%fakesection Commutative diagrams
\usepackage{diagrams}
\newarrow{Inj}C--->
\newarrow{Surj}----{>>}
\newarrow{Proj}----{>>}
\newarrow{Embed}>--->
\newarrow{Incl}C--->
\newarrow{Mapsto}|--->
\newarrow{Isom}<--->
\newarrow{Dotted}....>
\newarrow{Dashed}{}{dash}{}{dash}>
\newarrow{Line}-----
\usepackage{tikz-cd}
\usetikzlibrary{decorations.pathmorphing}
\tikzcdset{row sep/normal=3.14em,column sep/normal=3.14em}

%%fakesection Page layout
\usepackage[headsepline]{scrlayer-scrpage}
\renewcommand{\headfont}{}
\addtolength{\textheight}{3.14cm}
\setlength{\footskip}{0.5in}
\setlength{\headsep}{10pt}
\ihead{Evan Chen}
\automark{section}
\chead{}
\ohead{\headmark}
\cfoot{\pagemark}

%%fakesection Fancy section and chapter heads
\addtokomafont{chapterprefix}{\raggedleft}
\renewcommand*{\sectionformat}{\color{purple}\S\thesection\autodot\enskip}

\RedeclareSectionCommand[beforeskip=0.5em]{chapter}
\renewcommand*{\chapterformat}{%
\mbox{\scalebox{1.5}{\chapappifchapterprefix{\nobreakspace}}%
\scalebox{2.718}{\color{purple}\thechapter\autodot}\enskip}}

\addtokomafont{partprefix}{\rmfamily}
\renewcommand*{\partformat}{\color{purple}\scalebox{2.5}{\thepart}%
	\enlargethispage{2em}}
\newcommand{\problemhead}{Problems to think about}

%%fakesection Mini ToC
\usepackage[tight]{minitoc}
\noptcrule
\mtcsetfont{parttoc}{section}{\small\rmfamily\upshape\mdseries}
\renewcommand*{\partheadstartvskip}{}
\renewcommand*{\partheadendvskip}{}
\renewcommand\beforeparttoc{\hspace{\fill}\rule{0.95\linewidth}{2pt}\hspace{\fill}}
\let\oldpart\part
\renewcommand{\part}[1]{\oldpart{#1}\parttoc\cleardoublepage}

%%fakesection SNSD
\newcommand{\snsd}[2][]%
{\includegraphics[#1]{/home/evan/Pictures/TopologicalGG/#2}}

%%fakesection Theorems
\usepackage{thmtools}
\usepackage[framemethod=TikZ]{mdframed}

\theoremstyle{definition}
\mdfdefinestyle{mdbluebox}{%
	roundcorner = 10pt,
	linewidth=1pt,
	skipabove=12pt,
	innerbottommargin=9pt,
	skipbelow=2pt,
	nobreak=true,
	linecolor=blue,
	backgroundcolor=TealBlue!5,
}
\declaretheoremstyle[
	headfont=\sffamily\bfseries\color{MidnightBlue},
	mdframed={style=mdbluebox},
	headpunct={\\[3pt]},
	postheadspace={0pt}
]{thmbluebox}

\mdfdefinestyle{mdrecbox}{%
	linewidth=0.5pt,
	skipabove=12pt,
	frametitleaboveskip=5pt,
	frametitlebelowskip=0pt,
	skipbelow=2pt,
	frametitlefont=\bfseries,
	innertopmargin=4pt,
	innerbottommargin=8pt,
	nobreak=true,
	linecolor=RawSienna,
	backgroundcolor=Salmon!5,
}
\declaretheoremstyle[
	headfont=\bfseries\color{RawSienna},
	mdframed={style=mdrecbox},
	headpunct={\\[3pt]},
	postheadspace={0pt},
]{thmrecbox}

\declaretheorem[%
style=thmbluebox,name=Theorem,numberwithin=section]{theorem}
\declaretheorem[style=thmbluebox,name=Lemma,sibling=theorem]{lemma}
\declaretheorem[style=thmbluebox,name=Proposition,sibling=theorem]{proposition}
\declaretheorem[style=thmbluebox,name=Corollary,sibling=theorem]{corollary}
\declaretheorem[style=thmrecbox,name=Example,sibling=theorem]{example}

\mdfdefinestyle{mdinlinebox}{%
	skipabove=8pt,
	linewidth=3pt,
	rightline=false,
	leftline=true,
	topline=false,
	bottomline=false,
	linecolor=gray,
	backgroundcolor=black!2
}
\declaretheoremstyle[
	headfont=\bfseries,
	bodyfont=\normalfont\small,
	spaceabove=0pt,
	spacebelow=0pt,
	mdframed={style=mdinlinebox}
]{thminlinebox}
\theoremstyle{theorem}
\declaretheorem[name=Question,sibling=theorem,style=thminlinebox]{ques}
\declaretheorem[name=Exercise,sibling=theorem,style=thminlinebox]{exercise}

\theoremstyle{definition}
\newtheorem{claim}[theorem]{Claim}
\newtheorem{definition}[theorem]{Definition}
\newtheorem{fact}[theorem]{Fact}
\newtheorem{remark}[theorem]{Remark}
\newtheorem{abuse}[theorem]{Abuse of Notation}

\newtheorem{problem}{Problem}[chapter]
\renewcommand{\theproblem}{\thechapter\Alph{problem}}
\newtheorem{sproblem}[problem]{Problem}
\newtheorem{dproblem}[problem]{Problem}
\renewcommand{\thesproblem}{\theproblem$^{\star}$}
\renewcommand{\thedproblem}{\theproblem$^{\dagger}$}
\newcommand{\listhack}{$\empty$\vspace{-2em}}

%%fakesection Answers
\usepackage{answers}
\Newassociation{hint}{answeritem}{tex/backmatter/all-hints}
\Newassociation{sol}{answeritem}{tex/backmatter/all-solns}
\renewcommand{\solutionextension}{out}
\renewenvironment{answeritem}[1]{\item[\bf #1.]}{}

%%fakesection Hot chili peppers
\reversemarginpar
\newcommand{\prechili}{\vspace*{0.3em}\hspace*{1.5em}}
\newcommand{\nochili}{\hspace*{1.5em}}
\newcommand{\chili}{\includegraphics[width=1.5em]{media/chili.png}}
\newcommand{\gim}{\marginpar{\prechili\nochili\nochili\chili}}
\newcommand{\yod}{\marginpar{\prechili\nochili\chili\chili}}
\newcommand{\kurumi}{\marginpar{\prechili\chili\chili\chili}}

%%fakesection Table of contents
% First add ToC to ToC
\makeatletter
\usepackage{etoolbox}
\pretocmd{\tableofcontents}{%
  \if@openright\cleardoublepage\else\clearpage\fi
  \pdfbookmark[0]{\contentsname}{toc}%
}{}{}%
\makeatother
\setcounter{tocdepth}{1}
% Fix spacing issues in To
\usepackage[tocindentauto]{tocstyle}
\usetocstyle{KOMAlike}

%%fakesection Asymptote definitions
\usepackage{patch-asy}
\numberwithin{asy}{chapter}
\renewcommand{\theasy}{\thechapter\Alph{asy}}
\begin{asydef}
	import olympiad;
	import cse5;
	size(6cm);
	usepackage("amsmath");
	usepackage("amssymb");
	defaultpen(fontsize(11pt));
	pointfontsize=11;
	void bigbox(string L) {
		filldraw((-4,3)--(4,3)--(4,-3)--(-4,-3)--cycle,
			lightblue+opacity(0.2), black);
		label(L, (4,3), dir(-45));
	}
	void bigblob(string L) {
		filldraw((-4,3)..(-5,1)..(-4.7,0)..(-4.1,-2)..(-4,-3)
			..(-1,-3.3)..(1,-3.2)..(4,-3)
			..(4.3,-1)..(4.1,0.5)..(4,3)
			..(2,3.3)..(0,3.1)..(-3,2.9)..cycle,
			lightblue+opacity(0.2), black);
		label(L, (4,3), dir(5));
	}

	real[] default_ticks = {-4,-3,-2,-1,0,1,2.3,4};
	void cplane(real theta=-45, real[] T = default_ticks, real xmin = -4.5, real xmax = 4.5, real S = 3) {
		graph.xaxis("Re", xmin, xmax, graph.Ticks(format="%", Ticks=T, Size=S), Arrows);
		graph.yaxis("Im", xmin, xmax, graph.Ticks(format="%", Ticks=T, Size=S), Arrows);
		Drawing("0", origin, dir(theta));
		markscalefactor *= xmax;
	}

	pair opendot(pair P, pen p = defaultpen) {
		dot(P, p, UnFill);
		dot(P, p, Draw);
		return P;
	}
\end{asydef}
\def\asydir{asy}

%%fakesection Bibliography
\usepackage[backend=biber,style=alphabetic]{biblatex}
\DeclareLabelalphaTemplate{
	\labelelement{
		\field[final]{shorthand}
		\field{label}
		\field[strwidth=2,strside=left]{labelname}
	}
	\labelelement{
		\field[strwidth=2,strside=right]{year}    
	}
}
\DeclareFieldFormat{labelalpha}{\textbf{\scriptsize #1}}
\addbibresource{references.bib}
\addbibresource{images.bib}

%%fakesection Napkin Macros
\newcommand{\Zc}[1]{\ZZ_{#1}}
\newcommand{\Zm}[1]{\ZZ_{#1}^\times}
\newcommand{\pre}{^{\text{pre}}}
\newcommand{\normalin}{\trianglelefteq}
\newcommand{\triv}{\mathrm{triv}}
\newcommand{\largeotimes}[1]{\mathop{\otimes}\limits_{#1}}
\newcommand{\Mat}{\mathrm{Mat}}
\newcommand{\PGL}[1]{\mathbf{PGL}_{#1}(\CC)}

\DeclareMathOperator{\Stab}{Stab}
\DeclareMathOperator{\FixPt}{FixPt}
\DeclareMathOperator{\refl}{refl}
\DeclareMathOperator{\Fun}{Fun}
\DeclareMathOperator{\Irrep}{Irrep}
\DeclareMathOperator{\Res}{Res}
\DeclareMathOperator{\Reg}{Reg}
\DeclareMathOperator{\Classes}{Classes}
\newcommand{\FunCl}{\Fun_{\mathrm{class}}}
\newcommand{\Homrep}{\Hom_{\mathrm{rep}}}
\newcommand{\ab}{^\text{ab}} % abelianization
\DeclareMathOperator{\ev}{ev}
\DeclareMathOperator{\Wind}{\mathbf I}
\newcommand{\Ctriv}{\CC_{\mathrm{triv}}}
\newcommand{\Csign}{\CC_{\mathrm{sign}}}

% Alg geom macros
\newcommand{\VV}{\mathcal V}
\newcommand{\Vp}{\mathcal V_{\text{pr}}}
\newcommand{\Aff}{\mathbb A}
\newcommand{\II}{\mathcal I_{\text{rad}}}
\newcommand{\RP}{\mathbb{RP}}
\newcommand{\CP}{\mathbb{CP}}
\newcommand{\Cl}{\mathrm{Cl}}
\newcommand{\restrict}[1]{\bgroup\restriction_{#1}\egroup}
\DeclareMathOperator{\Opens}{OpenSets}
\DeclareMathOperator{\Proj}{Proj}
\DeclareMathOperator{\res}{res}
\newcommand{\mm}{\mathfrak m}
\newcommand{\sh}{^\mathrm{sh}}

\renewcommand{\Re}{\opname{Re}}
\renewcommand{\Im}{\opname{Im}}

%% Category Theory macros
\renewcommand{\AA}{\mathcal A}
\newcommand{\BB}{\mathcal B}
\newcommand{\obj}{\operatorname{obj}}
\newcommand{\op}{^{\mathrm{op}}}


\usepackage{tkz-graph}
\pgfarrowsdeclare{biggertip}{biggertip}{%
  \setlength{\arrowsize}{1pt}  
  \addtolength{\arrowsize}{.1\pgflinewidth}  
  \pgfarrowsrightextend{0}  
  \pgfarrowsleftextend{-5\arrowsize}  
}{%
  \setlength{\arrowsize}{1pt}  
  \addtolength{\arrowsize}{.1\pgflinewidth}  
  \pgfpathmoveto{\pgfpoint{-5\arrowsize}{4\arrowsize}}  
  \pgfpathlineto{\pgfpointorigin}  
  \pgfpathlineto{\pgfpoint{-5\arrowsize}{-4\arrowsize}}  
  \pgfusepathqstroke  
}  
\tikzset{
	EdgeStyle/.style = {>=biggertip}
}

\usepackage[all,cmtip,2cell]{xy}
\UseTwocells
\newcommand\nattfm[5]{\xymatrix@C+2pc{#1 \rtwocell<4>^{#2}_{#4}{\; #3} & #5}}

%% Alg Top macros
\DeclareMathOperator{\Cells}{Cells}
\newcommand{\HdR}{H_{\mathrm{dR}}}

%% Alg NT macros
\renewcommand{\aa}{\mathfrak a}
\newcommand{\bb}{\mathfrak b}
\newcommand{\pp}{\mathfrak p}
\newcommand{\qq}{\mathfrak q}

\newcommand{\Frob}{\mathrm{Frob}}
\DeclareMathOperator{\Norm}{N}
\DeclareMathOperator{\Ram}{Ram}
\newcommand{\TrK}{\Tr_{K/\QQ}}
\newcommand{\NK}{\Norm_{K/\QQ}}
\newcommand{\kf}{\mathfrak f}
\newcommand{\kP}{\mathfrak P}
\newcommand{\kQ}{\mathfrak Q}

%% Diff geo macros
\newcommand{\ee}{\mathbf e}
\newcommand{\fpartial}[2]{\frac{\partial #1}{\partial #2}}

%% ST macros
\newcommand{\CH}{\mathsf{CH}}
\newcommand{\ZFC}{\mathsf{ZFC}}

\newcommand{\Name}{\text{Name}}
\newcommand{\Po}{\mathbb P}
\newcommand{\nrank}{\opname{n-rank}} % ranks
\newcommand{\PP}{\mathcal P}

\newcommand{\EmptySet}{\mathrm{EmptySet}}
\newcommand{\PowerSet}{\mathrm{PowerSet}}
\newcommand{\Pairing}{\mathrm{Pairing}}
\newcommand{\Infinity}{\mathrm{Infinity}}
\newcommand{\Extensionality}{\mathrm{Extensionality}}
\newcommand{\Foundation}{\mathrm{Foundation}}
\newcommand{\Union}{\mathrm{Union}}
\newcommand{\Comprehension}{\mathrm{Comprehension}}
\newcommand{\Replacement}{\mathrm{Replacement}}

\usepackage{mathrsfs}

\newcommand\MM{\mathscr M}
\newcommand\llex{<_{\text{lex}}}

\DeclareMathOperator{\cof}{cof}

%% Quantum macros
\usepackage{braket}
\newcommand{\cvec}[1]{\begin{bmatrix} #1 \end{bmatrix}}
\newcommand{\pair}[2]{\begin{bmatrix} #1 \\ #2 \end{bmatrix}}
\newcommand{\zup}{\ket\uparrow}
\newcommand{\zdown}{\ket\downarrow}
\newcommand{\xup}{\ket\rightarrow}
\newcommand{\xdown}{\ket\leftarrow}
\newcommand{\yup}{\ket\otimes}
\newcommand{\ydown}{\ket\odot}
\newcommand{\UCNOT}{U_{\mathrm{CNOT}}}
\newcommand{\UQFT}{U_{\mathrm{QFT}}}

\def\asydir{}
\addbibresource{../../references.bib}
\renewcommand{\gim}{($\ast$)}

\begin{document}
\title{Singular Homology}
\maketitle

\chapter{Some Topological Constructions}
\section{Product Topology}
\section{Quotient Topology}
\section{Wedge Product}
Figure 8
\section{CW Complexes}

Torus, Klein, RP


\chapter{More on Computing Homology Groups}

\section{Excision}
relative homology groups, and excision

wedge products


\section{Degree}
\prototype{$z \mapsto z^d$ has degree $d$.}
For any $n > 0$ map $f$ from $S^n$ to itself
\[ f_\ast : \underbrace{H^n(S^n)}_{\cong \ZZ} \to \underbrace{H^n(S^n)}_{\cong \ZZ} \]
which must be multiplication by some constant $d$.
This $d$ is called the \vocab{degree} of $f$, denoted $\deg f$.
\begin{ques}
	Show that $\deg(f \circ g) = \deg(f) \deg(g)$.
\end{ques}

\begin{example}
	[Degree]
	\listhack
	\begin{enumerate}[(a)]
		\ii For $n=1$, the map $z \mapsto z^k$ (viewing $S^1 \subseteq \CC$)
		has degree $k$.
		\ii A reflection map $(x_0, x_1, \dots, x_n) \mapsto (-x_0, x_1, \dots, x_n)$
		has degree $-1$; we won't prove this, but geometrically this should be clear.
		\ii The antipodal map $x \mapsto -x$ has degree $(-1)^{n+1}$
		since it's the composition of $n+1$ reflections as above.
		We denote this map by $-\id$.
	\end{enumerate}
\end{example}

Obviously, if $f$ and $g$ are homotopic, then $\deg f = \deg g$.
In fact, a theorem of Hopf says that this is a classifying invariant:
anytime $\deg f = \deg g$, we have that $f$ and $g$ are homotopic.

One nice application of this:
\begin{theorem}
	[Hairy Ball Theorem]
	Whenever $n$ is even, $S^n$ doesn't have a continuous field
	of nonzero tangent vectors.
\end{theorem}
\begin{proof}
	If the vectors are nonzero then WLOG they have norm $1$;
	that is for every $x$ we have an orthogonal unit vector $v(x)$.
	Then we can construct a homotopy map $F : S^n \times [0,1] \to S^n$ by
	\[ (x,t) \mapsto (\cos \pi t)x + (\sin \pi t) v(x). \]
	which gives a homotopy from $\id$ to $-\id$.
	So $\deg \id = \deg(-\id)$, which means $1 = (-1)^{n+1}$.
\end{proof}

\section{Cellular Homology}
Before starting, we state the following lemma.
\begin{lemma}
	[CW Homology Groups]
	Let $X$ be a CW complex. Then
	\begin{align*}
		H_k(X^n, X^{n-1}) &\cong
		\begin{cases}
			\ZZ^{\oplus\text{\#$n$-cells of $X$}} & k = n \\
			0 & \text{otherwise}.
		\end{cases} \\
		\intertext{and}
		H_k(X^n) &\cong
		\begin{cases}
			H_k(X) & k \le n-1 \\
			0 & k \ge n+1.
		\end{cases}
	\end{align*}
\end{lemma}
\begin{proof}
	I'll prove just the case where $X$ is finite-dimensional for simplicity.
	The first part is immediate by noting that $(X^n, X^{n-1})$ is a good pair
	and $X^n/X^{n-1}$ is a wedge sum of two spheres.
	For the second part, fix $k$ and note that, as long as $n \le k-1$ or $n \ge k+2$,
	\[
		\underbrace{H_{k+1}(X^n, X^{n-1})}_{=0}
		\to H_k(X^{n-1})
		\to H_k(X^n)
		\to \underbrace{H_{k}(X^n, X^{n-1})}_{=0}.
	\]
	So we have isomorphisms
	\[ H_k(X^{k-1}) \cong H_k(X^{k-2}) \cong \dots \cong H_k(X^0) = 0 \]
	and
	\[ H_k(X^{k+1}) \cong H_k(X^{k+2}) \cong \dots \cong H_k(X). \qedhere \]
\end{proof}

So, we know that the groups $H_k(X^n, X^{n-1})$ are super nice:
they are free abelian with basis given by the cells of $X$.

In light of this, let's look at our long exact sequence and try to
string together maps between these.
Consider the following diagram.

\begin{diagram}
	\small
	&& \underbrace{H_3(X^2)}_{=0} &&&&&&&& \\
	&& \dTo~0 &&&&&&&& \\
	\boxed{H_4(X^4, X^3)} & \rTo^{\partial_4} & H_3(X^3) & \rSurj
		& \underbrace{H_3(X^4)}_{\cong H_3(X)} &
		\rTo~0 & \underbrace{H_3(X^4, X^3)}_{= 0} &&&& \\
	& \rdDotted~{d_4} & \dTo~0 &&&&&&&& \\
	&& \boxed{H_3(X^3, X^2)} &&&& \underbrace{H_1(X^0)}_{=0} &&&& \\
	&& \dTo^{\partial_3} & \rdDotted~{d_3} &&& \dTo~0 &&&& \\
	\underbrace{H_2(X^1)}_{=0} & \rTo~0 & H_2(X^2) & \rInj &
		\boxed{H_2(X^2, X^1)} & \rTo^{\partial_2} & H_1(X^1) & \rSurj &
		\underbrace{H_1(X^2)}_{\cong H_1(X)} & \rTo~0 &
		\underbrace{H_1(X^2, X^1)}_{=0} \\
	&& \dSurj &&& \rdDotted~{d_2} & \dInj &&&& \\
	&& \underbrace{H_2(X^3)}_{\cong H_2(X)} &&&& \boxed{H_1(X^1, X^0)} &&&& \\
	&& \dTo~0 &&&& \dTo^{\partial_1} & \rdDotted~{d_1} &&& \\
	&& \underbrace{H_2(X^3, X^2)}_{=0} && \underbrace{H_0(\varnothing)}_{=0}
		& \rTo~0 & H_0(X^0) & \rInj & \boxed{H_0(X^0, \varnothing)} & \rTo^{\partial_0} & \dots \\
	&&&&&& \dSurj &&&& \\
	&&&&&& \underbrace{H_0(X^1)}_{\cong H_0(X)} &&&& \\
	&&&&&& \dTo~0 &&&& \\
	&&&&&& \underbrace{H_0(X^1, X^0)}_{=0} &&&& \\
\end{diagram}
Here we have $X^{-1} = \varnothing$.
The idea is that we have taken all the exact sequences generated by adjacent
skeletons, and strung them together at the groups $H_k(X^k)$,
with half the exact sequences being laid out vertically
and the other half horizontally.

In that case, composition generates a sequence of dotted maps
between the $H_k(X^k, X^{k-1})$ as shown.
\begin{ques}
	Show that the composition of two adjacent dotted arrows is zero.
\end{ques}

So the sequence of arrows $\dots \to H_4(X^4, X^3) \to H_3(X^3, X^2) \to \dots$
that we obtain is a chain complex, called the \vocab{cellular chain complex};
as mentioned before before all the homology groups are free,
but these ones are especially nice because for most reasonable CW complexes,
they are also finitely generated
(unlike the massive $C_\bullet(X)$ that we had earlier).
In other words, the $H_k(X^k, X^{k-1})$ are especially nice ``concrete'' free groups
that one can actually work with.

The other reason we care is that in fact:
\begin{theorem}[Cellular Chain Complex Gives $H_n(X)$]
	\label{thm:cellular_chase}
	The $k$th homology group of the cellular chain complex
	is isomorphic to $H_k(X)$.
\end{theorem}
\begin{proof}
	Follows from the diagram; \Cref{prob:diagram_chase}.
\end{proof}

In fact, one can describe explicitly what the maps $d_n$ are.
Recalling that $H_k(X^k, X^{k-1})$ has a basis the $k$-cells of $X$, we
can write the following formula.
\begin{theorem}
	[Cellular Boundary Formula]
	Let $e$ be a $k$-cell, and let $f_\beta$ denote all $(k-1)$-cells of $X$.
	Define $d_\beta$ to be the degree of the composed map
	\[ \partial D^n \xrightarrow{\text{attach}} X^{n-1} \surjto S_\beta^{n-1} \]
	where the first arrow is the attaching map for $e$
	and the second arrow is the quotient of collapsing $X^{n-1} - f_\beta$ to a point.
	Then \[ d_n : e \mapsto \sum_\beta d_\beta f_\beta \]
\end{theorem}

CP

Torus, Klein, RP2

\section\problemhead

\begin{problem}
	Show that a non-surjective map $f : S^n \to S^n$ has degree zero.
	\begin{hint}
		The space $S^n - \{x_0\}$ is contractible.
	\end{hint}
\end{problem}

\begin{problem}
	\label{prob:diagram_chase}
	Prove \Cref{thm:cellular_chase},
	showing that the homology groups of $X$
	coincide with the homology groups of the cellular chain complex.
	\begin{hint}
		You won't need to refer to any elements.
		Start with $H_2(X) \cong H_2(X^3) \cong
			H_2(X^2) / \ker \left[ H_2(X^2) \surjto H_2(X^3) \right]$, say.
		Take note of the marked injective and surjective arrows.
	\end{hint}
	\begin{sol}
		For concreteness, let's just look at the homology at $H_2(X^2, X^1)$
		and show it's isomorphic to $H_2(X)$.
		According to the diagram
		\begin{align*}
			H_2(X) &\cong H_2(X^3) \\
			&\cong H_2(X^2) / \ker \left[ H_2(X^2) \surjto H_2(X^3) \right] \\
			&\cong H_2(X^2) / \img \partial_3 \\
			&\cong \img\left[ H_2(X^2) \injto H_2(X^2, X^1) \right] / \img \partial_3 \\
			&\cong \ker(\partial_2) / \img\partial_3 \\
			&\cong \ker d_2 / \img d_3. \qedhere
		\end{align*}
	\end{sol}
\end{problem}

\begin{problem}[Moore Spaces]
	Let $G_1$, $G_2$, \dots, $G_N$ be a sequence of finitely generated abelian groups.
	Construct a space $X$ such that
	\[
		\wt H_n(X)
		\cong
		\begin{cases}
			G_n & 1 \le n \le N \\
			0 & \text{otherwise}.
		\end{cases}
	\]
\end{problem}

\chapter{Singular Cohomology}
Here's the main way to motivate this chapter.
Let $X = \CP^2$ and $Y = S^2 \wedge S^4$.
\begin{exercise}
	Show that $H_n(X) \cong H_n(Y)$ for every $n$.
\end{exercise}
However, as it turns out $X$ and $Y$ are not homotopy equivalent --- so how can we tell them apart?

In this chapter, we'll define a \emph{cohomology group} $H^n(X)$ and $H^n(Y)$.
We will still have $H^n(X) \cong H^n(Y)$ for every $X$ and $Y$:
in fact, the $H^n$'s are completely determined by the $H_n$'s
by the so-called \emph{Universal Coefficient Theorem}
However, it turns out that one can take all the cohomology groups and put
them together to form a \emph{cohomology ring} $H^\bullet$.
We will then see that $H^\bullet(X) \not\cong H^\bullet(Y)$ as rings.

\section{Cochain Complexes}
\begin{definition}
A \vocab{cochain complex} $A^\bullet$ is algebraically the same as a chain complex:
it is a sequence of abelian groups
\[ \dots \taking{\delta} A^{n-1} \taking\delta A^n \taking\delta A^{n+1} \taking\delta \dots. \]
such that $\delta^2 = 0$.
Notation-wise, we're now using subscripts, and use $\delta$ rather $\partial$.
We can still define the homology groups by
\[ H_n(A^\bullet) = \ker\left( A^n \taking\delta A^{n+1} \right)
	/ \img\left( A^{n-1} \taking\delta A^n \right). \]
\end{definition}

\textbf{The most common way to get a cochain complex
is to \emph{dualize} a chain complex.}
Specifically, pick an abelian group $G$;
note that $\Hom(-, G)$ is a contravariant functor,
and thus takes every chain complex
\[ \dots \taking\partial A_{n+1} \taking\partial
	A_n \taking\partial A_{n-1} \taking\partial \dots \]
into a cochain complex: letting $A^n = \Hom(A_n, G)$ we obtain
\[ \dots \taking\delta A^{n-1} \taking\delta
	A^n \taking\delta A^{n+1} \taking\delta \dots. \]
where $\delta(A_n \taking{f} G) = A_{n+1} \taking\partial A \taking{f} G$.
\begin{definition}
	Given a chain complex $A_\bullet$ of abelian groups and a group $G$,
	we define the \vocab{cohomology groups} $H^n(A_\bullet; G)$ to be the
	homology groups of the dual $A^\bullet$ obtained by applying $\Hom(-,G)$.
\end{definition}
This is admittedly horrible naming:
the \emph{cohomology groups of a chain complex}
are the \emph{homology groups of the cochain complex}
dual to that chain complex.

\section{Cohomology of Spaces}
\prototype{$C^0(X;G)$ all functions $X \to G$ while $H^0(X)$ are those functions $X \to G$
constant on path components.}

The case of interest is our usual geometric situation, with $C_\bullet(X)$.
\begin{definition}
	For a space $X$ and abelian group $G$,
	we define $C^\bullet(X;G)$ to be the dual complex to $C_\bullet(X)$,
	called the \vocab{singular cochain complex} of $X$;
	its elements are called \vocab{cochains}.

	Then we define the \vocab{cohomology groups}
	of the space $X$ as 
	\[ H^n(X; G) \defeq H^n(C_\bullet(X); G) = H_n(C^\bullet(X;G)). \]
\end{definition}

\begin{example}
	[$C^0(X; G)$, $C^1(X; G)$, and $H^0(X;G)$]
	Of course the archetypal example is to dualize the singular chain complex
	of a space $X$. In this situation, elements of $C^n$ are called \emph{cochains}.
	\begin{itemize}
		\ii $C_0(X)$ is the free abelian group on $X$,
		so $C^0(X) = \Hom(C_0(X), G)$ is a function that
		takes every point of $X$ to an element of $G$.
		\ii $C_1(X)$ is the free abelian group on $1$-simplices in $X$.
		So $C^1(X)$ needs to take every $1$-simplex to an element of $G$.
	\end{itemize}
	Let's now try to understand $\delta : C^1(X) \to C^0(X)$.

	Now given a cochain $\phi \in C^0(X)$, i.e.\ a homomorphism $\phi : C^0(X) \to G$,
	what is $\delta\phi : C^1(X) \to G$?
	Answer: 
	\[ \delta\phi : [v_0, v_1] \mapsto \phi([v_0]) - \phi([v_1]). \]
	Hence, elements of $\ker(C^0 \taking\partial C^1)$ are those cochains
	that are \emph{constant on path-connected components}.
\end{example}
In particular, much like $H_0(X)$, we have \[ H^0(X) \cong G^{\oplus r} \]
if $X$ has $r$ path-connected components (where $r$ is finite\footnote{%
	Something funny happens if $X$ has \emph{infinitely} many path-connected components:
	say $X = \bigsqcup_\alpha X_\alpha$ over an infinite indexing set.
	In this case we have
	$H_0(X) = \bigoplus_\alpha G$ while $H^0(X) = \prod_\alpha G$.
	For homology we get a \emph{direct sum} while
	for cohomology we get a \emph{direct product}.

	These are actually different for infinite indexing sets.
	For general modules $\bigoplus_\alpha M_\alpha$ is \emph{defined} to only allow
	to have \emph{finitely many} zero terms.
	(This was never mentioned earlier in the Napkin,
	since I only ever defined $M \oplus N$ and extended it to finite direct sums.)
	No such restriction holds for $\prod_\alpha G_\alpha$ a product of groups.
	This corresponds to the fact that $C_0(X)$ is formal linear sums of $0$-chains
	(which, like all formal sums, are finite)
	from the path-connected components of $G$.
	But a cochain of $C^0(X)$ is a \emph{function}
	from each path-connected component of $X$ to $G$,
	where there is no restriction.
}).

To the best of my knowledge, the higher cohomology groups $H^n(X; G)$
(or even the cochain groups $C^n(X; G) = \Hom(C_n(X), G)$) are harder to describe concretely.

\begin{abuse}
	In this chapter the only cochain complexes
	we will consider are dual complexes as above.
	So, any time we write a chain complex $A^\bullet$ it is implicitly given
	by applying $\Hom(-,G)$ to $A_\bullet$.
\end{abuse}

\section{Cohomology of Spaces is Functorial}
We now check that the cohomology groups still exhibit the same nice functorial behavior.
First, let's categorize the previous results we had:

\begin{ques}
	Define $\catname{CoCmplx}$
	the category of cochain complexes.
\end{ques}

\begin{exercise}
	Interpret $\Hom(-,G)$ as a contravariant functor
	from \[ \Hom(-,G) : \catname{Cmplx}\op \to \catname{CoCmplx}. \]
	This means in particular that given a chain map $f : A_\bullet \to B_\bullet$,
	we naturally obtain a dual map $f^\vee : B^\bullet \to A^\bullet$.
\end{exercise}

\begin{ques}
	Interpret $H^n : \catname{CoCmplx} \to \catname{Grp}$ as a functor.
	Compose these to get a contravariant functor
	$H^n(-;G) : \catname{Cmplx}\op \to \catname{Grp}$.
\end{ques}

Then in exact analog to our result that $H_n : \catname{hTop} \to \catname{Grp}$ we have:
\begin{theorem}[$H^n (-;G): \catname{hTop}\op \to \catname{Grp}$]
	For every $n$, $H^n(-;G)$ is a contravariant functor
	from $\catname{hTop}\op$ to $\catname{Grp}$.
\end{theorem}
\begin{proof}
	The idea is to leverage the work we already did in constructing
	the prism operator earlier.
	First, we construct the entire sequence of functors
	from $\catname{Top}\op \to \catname{Grp}$:
	\begin{diagram}
		\catname{Top}\op & \rTo^{C_\bullet} & \catname{Cmplx}\op & \rTo^{\Hom(-;G)}
		& \catname{CoCmplx} & \rTo^{H^n} & \catname{Grp} \\
		X && C_\bullet(X) && C^\bullet(X;G) && H^n(X;G) \\
		\dTo^f &\rMapsto& \dTo^{f_\sharp} &\rMapsto&
		\uTo^{f^\sharp} &\rMapsto& \uTo^{f^\ast} \\
		Y && C_\bullet(Y) && C^\bullet(Y;G) && H^n(Y;G).
	\end{diagram}
	Here $f^\sharp = (f_\sharp)^\vee$, and $f^\ast$
	is the resulting induced map on homology groups of the cochain complex.

	So as before all we have to show is that $f \simeq g$,
	then $f^\ast = g^\ast$.
	Recall now that there is a prism operator such that
	$f_\sharp - g_\sharp = P \partial + \partial P$.
	If we apply the entire functor $\Hom(-;G)$ we get that
	$f^\sharp - g^\sharp = \delta P^\vee + P^\vee \delta$
	where $P^\vee : C^{n+1}(Y;G) \to C^n(X;G)$.
	So $f^\sharp$ and $g^\sharp$ are chain homotopic thus $f^\ast = g^\ast$.
\end{proof}


\section{Universal Coefficient Theorem}
We now wish to show that the cohomology groups are determined up to isomorphism
by the homology groups: given $H_n(A_\bullet)$, we can extract $H^n(A_\bullet; G)$.
This is achieved by the \emph{Universal Coefficient Theorem}.
\begin{theorem}
	[Universal Coefficient Theorem]
	Let $A_\bullet$ be a chain complex of \emph{free} abelian groups,
	and let $G$ be another abelian group.
	Then there is a natural short exact sequence
	\[
		0 \to \Ext(H_{n-1}(A_\bullet), G) \to H^n(A_\bullet; G)
		\taking{h} \Hom(H_n(A_\bullet), G) \to 0. \]
	In addition, this exact sequence is \emph{split}
	so in particular
	\[ H^n(C_\bullet; G) \cong \Ext(H_{n-1}(A_\bullet, G)
		\oplus \Hom(H_n(A_\bullet), G). \]
\end{theorem}
Fortunately, in our case of interest, $A_\bullet$ is $C_\bullet(X)$
which is by definition free.

There are two things we need to explain, what the map $h$ is and the map $\Ext$ is.

It's not too hard to guess how \[ h : H^n(A_\bullet; G) \to \Hom(H_n(A_\bullet), G) \] is defined.
An element of $H^n(A_\bullet;G)$ is represented by a function which sends a cycle
in $A_n$ to an element of $G$.
The content of the theorem is to show that $h$ is surjective with kernel $\Ext(H_{n-1}(A_\bullet), G)$.

What about $\Ext$?
It turns out that $\Ext(-,G)$ is the so-called \vocab{Ext functor}, defined as follows.
Let $H$ be an abelian group, and consider a \vocab{free resolution} of $H$,
by which we mean an exact sequence
\[ \dots \taking{f_2} F_1 \taking{f_1} F_0 \taking{f_0} H \to 0. \]
Then we can apply $\Hom(-,G)$ to get a cochain complex
\[ \dots \xleftarrow{f_2^\vee} \Hom(F_1, G) \xleftarrow{f_1}
	\Hom(F_0, G) \xleftarrow{f_0} \Hom(H,G) \to 0. \]
but \emph{this cochain complex need not be exact}
(in categorical terms, $\Hom(-,G)$ does not preserve exactness).
We define \[ \Ext(H,G) \defeq \ker(f_2^\vee) / \img(f_1^\vee) \]
and it's a theorem that this doesn't depend on the choice of the free resolution.
There's a lot of homological algebra that goes into this,
which I won't take the time to discuss;
but the upshot of the little bit that I did include is that the $\Ext$
functor is very easy to compute in practice, since
you can pick any free resolution you want and compute the above.

%By ``natural'', we mean that if $f : A_\bullet \to B_\bullet$ is a chain map,
%then we obtain a commutative diagram
%\begin{diagram}
%	0 & \rTo & \Ext(H_{n-1}(A_\bullet), G) & \rTo
%		& H^n(A_\bullet;G) & \rTo & \Hom(H_n(A_\bullet), G) & \rTo & 0 \\
%	& & \uTo^{ \Ext(f_\ast, G) } & & \uTo^{f^\ast} & & \uTo^{\Hom(f_\ast, G)} & & \\
%	0 & \rTo & \Ext(H_{n-1}(B_\bullet), G) & \rTo
%		& H^n(A_\bullet;G) & \rTo & \Hom(H_n(B_\bullet), G) & \rTo & 0 \\
%\end{diagram}
%where $f_\ast$ is the induced arrow $H_n(A_\bullet) \to H_n(B_\bullet)$.

\begin{lemma}
	[Computing the $\Ext$ Functor]
	For any abelian groups $G$, $H$, $H'$ we have
	\begin{enumerate}[(a)]
		\ii $\Ext(H \oplus H, G) = \Ext(H, G) \oplus \Ext(H', G)$.
		\ii $\Ext(H,G) = 0$ for $H$ free, and
		\ii $\Ext(\Zc n, G) = G / nG$.
	\end{enumerate}
\end{lemma}
\begin{proof}
	For (a), note that if $\dots \to F_1 \to F_0 \to H \to 0$
	and $\dots \to F_1' \to F_0' \to F_0' \to H' \to 0$ are free resolutions,
	then so is $F_1 \oplus F_1' \to F_0 \oplus F_0' \to H \oplus H' \to 0$.

	For (b), note that $0 \to H \to H \to 0$ is a free resolution.
	
	Part (c) follows by taking the free resolution
	\[ 0 \to \ZZ \taking{\times n} \ZZ \to \Zc n \to 0 \]
	and applying $\Hom(-,G)$ to it.
	\begin{ques}
		Finish the proof of (c) from here. \qedhere
	\end{ques}
\end{proof}

\begin{ques}
	Some $\Ext$ practice: compute
	$\Ext(\ZZ^{\oplus 2015}, G)$ and $\Ext(\Zc{30}, \Zc 4)$.
\end{ques}

\section{Example Computation of Cohomology Groups}
\prototype{Possibly $H^n(S^m)$.}
Before we start doing some computations, we take the time
to define reduced cohomology groups as well.
For $X$ nonempty we define the \vocab{reduced cohomology groups} by $\wt H^n(X;G)$
by $H^n(\wt C_\bullet(X); G)$.
To be concrete, the augmented singular chain complex is
\[ \dots \taking\partial C_1(X) \taking\partial C_0(X) \taking\eps \ZZ \to 0 \]
dualizes to
\[
	\dots \xleftarrow\delta C^1(X;G) \xleftarrow\delta C^0(X;G)
	\xleftarrow{\eps^\vee} \underbrace{\Hom(\ZZ, G)}_{\cong G}
	\leftarrow 0.
\]
Since the $\ZZ$ we add is also free,
the Universal Coefficient Theorem still applies.

\begin{ques}
	For each $n \ge 1$, $\wt H^n(X;G) \cong H^n(X;G)$.
\end{ques}

As for $\wt H^0(X;G)$ we use the Universal Coefficient Theorem;
since the homology at $\ZZ$ is trivial, the $\Ext$ term vanishes and
$\wt H^0(X;G) \cong \Hom(\wt H_0(X), G)$.
Geometrically, $\wt H^0(X;G)$ consists of functions $X \to G$
which are constant on path-connected components,
modulo the functions which are constant on all of $X$.

\begin{example}
	[Cohomolgy Groups of $S^m$]
	It is straightforward to compute $H^n(S^m)$ now:
	all the $\Ext$ terms vanish since $H_n(S^m)$ is always free,
	and hence we obtain that 
	\[ H^n(S^m) \cong \Hom(H_n(S^m), G) \cong
		\begin{cases}
			G & n=m, n=0 \\
			0 & \text{otherwise}.
		\end{cases}
	\]
	By UCT for reduced groups, we also have
	\[ \wt H^n(S^m) \cong \Hom(\wt H_n(S^m), G) \cong
		\begin{cases}
			G & n=m \\
			0 & \text{otherwise}.
		\end{cases}
	\]
	since $\Hom(\ZZ, G)$.
\end{example}

\begin{example}
	[Cohomolgy Groups of Torus]
	This example has no nonzero $\Ext$ terms either,
	since this time $H^n(S^1 \times S^1)$ is always free.
	So we obtain
	\[ H^n(S^1 \times S^1) \cong \Hom(H_n(S^1 \times S^1), G)
		\cong
		\begin{cases}
			G & n = 0,2 \\
			G^{\oplus 2} & n = 1.
		\end{cases}
	\]
\end{example}

\begin{example}
	[Cohomolgy Groups of Klein Bottle]
	This example will actually have $\Ext$ term.
	Recall that if $K$ is a Klein Bottle then its homology groups are
	$\ZZ$ in dimension $n=0$ and $\ZZ \oplus \Zc 2$ in $n=1$, and $0$ elsewhere.

	For $n=0$, we again just have $H^0(K;G) \cong \Hom(\ZZ, G) \cong G$.
	For $n=1$, the $\Ext$ term is $\Ext(H_0(K), G) \cong \Ext(\ZZ, G) = 0$
	so \[ H^1(K;G) \cong \Hom(\ZZ \oplus \Zc2, G) \cong G \oplus \Hom(\Zc2, G). \]
	We have that $\Hom(\Zc2,G)$ is the subgroup
	of elements of order $2$ in $G$ (and $0 \in G$).

	But for $n=2$, we have our first interesting $\Ext$ group:
	the exact sequence is
	\[ 0 \to \Ext(\ZZ \oplus \Zc 2, G) \to H^2(X;G) \to \underbrace{H_2(X)}_{=0} \to 0. \]
	Thus, we have
	\[ H^2(X;G) \cong \left( \Ext(\ZZ,G) \oplus \Ext(\Zc2,G) \right) \oplus 0
		\cong G/2G. \]
	All the higher groups vanish.
	In summary:
	\[
		H^n(X;G) \cong
		\begin{cases}
			G & n = 0 \\
			G \oplus \Hom(\Zc2, G) & n = 1 \\
			G/2G & n = 2 \\
			0 & n \ge 3. 
		\end{cases}
	\]
\end{example}

From the previous example we can extract the following observation.
\begin{lemma}
	[$0$th and $1$st Homology Groups are just Duals]
	For $n = 0$ and $n = 1$, we have
	\begin{align*}
		H^n(X;G) &\cong \Hom(H_n(X), G) \\
		\wt H^n(X;G) &\cong \Hom(\wt H_n(X), G).
	\end{align*}
\end{lemma}
\begin{proof}
	It's already been shown for $n=0$.
	For $n=1$, notice that $H_0(X)$ and $\wt H_0(X)$ are free,
	so the $\Ext$ term vanishes.
\end{proof}

\section\problemhead
\begin{sproblem}
	[Wedge Product Cohomology]
	For any $G$ and $n$ we have
	\[
		\wt H^n(X \vee Y; G)
		\cong
		\wt H^n(X; G) \oplus \wt H^n(Y; G).
	\]	
\end{sproblem}
\begin{sproblem}[$\Zc2$-Cohomology of $\RP^n$]
	Show that
	\[
		H^m(\RP^n, \Zc2)
		\cong
		\begin{cases}
			\ZZ & \text{$m=0$, or $m$ is odd and $m=n$} \\
			\Zc2 & \text{$0 < m < n$ and $m$ is odd} \\
			0 & \text{otherwise}.
		\end{cases}
	\]
\end{sproblem}

\chapter{The Cup Product}
I won't do much more in this chapter other than
introduce the cup product and say a few examples.
For an actual treatise, see \cite{ref:hatcher} or \cite{ref:maxim752}.

\section{Graded Rings}
\prototype{Polynomial rings are graded rings.}
In this chapter, the ring $R$ is commutative with identity;
but the remaining rings need not be commutative.
However, the word ``ring'' refers to a ring with identity,
or else we will use \vocab{pseudo-ring} to denote a ring that need not possess $1$.

\todo{define graded ring}

\todo{graded commutative}

\section{Cup Products}
Suppose now we replace our abelian group $G$ by a commutative ring $R$ with identity.
Then this gives us a way to multiply two cochains, as follows.
Suppose $\phi \in C^k(X;R)$ and $\psi \in C^\ell(X;R)$.
Then we can define their \vocab{cup product}
$\phi\smile\psi \in C^{k+\ell}(X;R)$ to be
\[
	(\phi\smile\psi)([v_0, \dots, v_{k+\ell}])
	= 
	\phi\left( [v_0, \dots, v_k] \right)
	\cdot
	\psi\left( [v_k, \dots, v_{k+\ell}] \right)
\]
where the multiplication is in $R$.

\begin{ques}
	Which $0$-cochain is the identity for $\smile$?
\end{ques}

This behaves well with respect to $\delta$:
\begin{lemma}[$\delta$ with Cup Products]
	We have
	$\delta(\phi\smile\psi) = \delta\phi\smile\psi + (-1)^k\phi\smile\delta\psi$.
\end{lemma}
\begin{proof}
	Direct $\sum$ computations.
\end{proof}
From this, it is not hard to see that we get an induced map
\[ H^k(X;R) \times H^\ell(X;R) \taking\smile H^{k+\ell}(X;R).  \]
We then define the \vocab{cohomology ring} $H^\bullet(X;R)$
whose elements are finite sums in 
\[ \bigoplus_{k \ge 0} H^k(X;R) \]
and with multiplication given by $\smile$.

Observe that $H^\bullet(X;R)$ is a \emph{graded ring}.
We say that an element $\alpha \in H^k(X;R)$ has dimension $k$,
and write $|\alpha|=k$.
It turns out (though we won't prove it; see \cite{ref:hatcher})
that the cup product is in fact graded commutative.

\begin{example}[Cohomology Ring of $S^n$]
	Consider $S^n$ for $n \ge 1$.
	The nontrivial cohomology groups are given by
	$H^0(S^n; \ZZ) \cong H^n(S^n, \ZZ) \cong \ZZ$.
	So as an abelian group
	\[ H^\bullet(S^n, \ZZ) \cong \ZZ \oplus \alpha \ZZ \]
	where $\alpha$ is the generator of $H^n(S^n, \ZZ)$.
	
	Now, observe that $|\alpha\smile\alpha| = 2n$, but
	since $H^{2n}(S^n, \ZZ) = 0$ we must have $\alpha\smile\alpha=0$.
	So even more succinctly,
	\[ H^\bullet(S^n, \ZZ) \cong \ZZ[\alpha]/(\alpha^2). \]
\end{example}

\begin{example}[Real and Complex Projective Space]
	Though we won't prove it, it turns out that
	\begin{align*}
		H^\bullet(\RP^n, \Zc2) \cong \Zc2[\alpha]/(\alpha^{n+1}) \\
		H^\bullet(\CP^n, \ZZ) \cong \ZZ[\beta]/(\beta^{n+1})
	\end{align*}
	where $|\alpha|=1$ is a generator of $H^1(\RP^n; \Zc2)$
	and $|\beta|=2$ is a generator of $H^2(\CP^n; \ZZ)$.
\end{example}

\section{Cup Product Structure}
In general, it's not so easy to figure out what the cup product looks like
for an arbitrary space.


\section{Kunneth Formula}
\section{The Borsuk-Ulam Theorem}

\section\problemhead

\end{document}
